\begin{helper}
In the fixed history cell, assume \(u_R\le u_B\).  With
\[
a=f(\ell_R,\ell_B)-10\varepsilon_1 k^2,
\]
there are natural numbers \(N_R,N_B\) such that
\[
N_R\le 3\varepsilon_1 k^2,\qquad N_B\le k^2,
\]
and the conditional probability of the fixed-count target is at most
\[
\binom{N_R}{u_B}p^{u_B}\binom{N_B}{u_R}p^{u_R}
(1-p)^{\max(0,a-u_B-u_R)}.
\]
Here blue is the later-color support controlled by the exceptional sets, and
red is the first-color support controlled by the \(k\)-set projection
universe.
\end{helper}
\begin{proof}
This is the color-reversed version of
\(\noderef{FixedSetHistoryCellFixedCountRedWeightSum}\), but we spell out the
objects to fix the bookkeeping.  Let \(H\) be the history cell and let
\(T_{u_R,u_B}\) be the same fixed-count target event.  Put
\[
L=\max(0,a-u_B-u_R).
\]

Use \(\noderef{FixedSetHistoryCellGlobalBlueExceptionalWitness}\) to replace
the branch-dependent blue exceptional sets by one global bound
\(N_R\le3\varepsilon_1k^2\).  For each complete red terminal branch, every
regular independent completion has its blue present staged-open support
inside a blue exceptional set of cardinality at most \(N_R\).  Therefore the
blue later support has at most \(\binom{N_R}{u_B}\) possible global labels.

The first-color red present support lies in the projected-pair universe
coming from \(I\).  Since \(|I|=k\), that universe has size at most \(k^2\);
write this size as \(N_B\).  Thus there are at most
\(\binom{N_B}{u_R}\) red first-color support labels.

For fixed support labels, partition the event by complete red terminal traces
and by prefix transcripts of the canonical absence-selection algorithm.  The
complete trace partition is supplied by
\(\noderef{FixedSetHistoryCellAdaptiveBranchPartition}\), and the selected
absent set \(Z_\tau\) is supplied by
\(\noderef{FixedSetHistoryCellCanonicalAbsenceSelection}\).  It is disjoint
from the two present supports and has size at least \(L\).

Each cell fixes a set \(P_\tau\) of present terminal coordinates and a
disjoint set \(A_\tau\) of absent terminal coordinates.  These sets include
all complete-trace values and all prefix-transcript values, not only the two
support labels and selected absences.  Applying
\(\noderef{FixedSetHistoryCellFixedCylinderCharge}\) to the terminal product
law bounds the raw cell mass by
\[
\Pr(H)\,p^{|P_\tau|}(1-p)^{|A_\tau|}.
\]

When the cell weights are summed for fixed support labels, the auxiliary
terminal coordinates fixed only to determine the trace or transcript are
reindexed over a finite residual set.  Their total contribution is
\[
\sum_W p^{|W|}(1-p)^{|W^c|}=(p+(1-p))^{|W\cup W^c|}=1.
\]
The mandatory blue and red present supports contribute \(p^{u_B}p^{u_R}\),
and the selected absences contribute at most \((1-p)^L\).  Hence the total
weight for fixed support labels is bounded by
\[
p^{u_B}p^{u_R}(1-p)^L.
\]

Summing over at most \(\binom{N_R}{u_B}\) blue support labels and at most
\(\binom{N_B}{u_R}\) red support labels gives the displayed total weight.
The cover, per-cell estimates, and total weight estimate are exactly the
hypotheses of \(\noderef{FixedSetHistoryCellBranchTranscriptSummation}\).
This bounds the raw probability of \(H\cap T_{u_R,u_B}\).  If \(\Pr(H)=0\),
the conditional probability is zero by
\(\noderef{TwoBiteConditionalProbability}\); otherwise divide by \(\Pr(H)\).
\end{proof}
